\chapter{Complexity — ARM Loops}
%%% generated by the blueprint pipeline from TCSlib.Complexity.SpaceComplexity.Machines.ARMLoops
\section{Overview}

Declarations of the Lean module \texttt{TCSlib.Complexity.SpaceComplexity.Machines.ARMLoops}.

\section{Declarations}

\begin{definition}[Bounded admissible configurations]
\lean{Complexity.LogProg.GuardLe}
\label{Complexity.LogProg.GuardLe}
\leanok
\uses{Complexity.LogProg.AConf, Complexity.LogProg.ARM, Complexity.LogProg.PreS}
Given a program with $m$ registers, an input bit string $x$, and a natural number $B$,
call a configuration admissible with bound $B$ if it satisfies the program’s
precondition for $x$ and every register contains a number at most $B$.
\end{definition}

\begin{lemma}[Register bound guard]
\lean{Complexity.LogProg.GuardLe.of_preS}
\label{Complexity.LogProg.GuardLe.of_preS}
\leanok
\uses{Complexity.LogProg.AConf, Complexity.LogProg.GuardLe, Complexity.LogProg.PreS}
\difficulty{1}
Let $A$ be a register program, $x$ an input, $B$ a natural number, and $a$ a
configuration consisting of a control state, register values, and an optional answer. If
the precondition of $A$ holds at $a$ for input $x$, and every register value is at most
$B$, then $a$ satisfies the guard requiring both conditions.
\end{lemma}

\begin{lemma}[Reachability under endpoint equality]
\lean{Complexity.LogProg.reach_of_eq}
\label{Complexity.LogProg.reach_of_eq}
\leanok
\uses{Complexity.LogProg.AConf, Complexity.LogProg.AReach}
\difficulty{1}
Let a deterministic register machine, oracle, input, and configuration property $G$ be
fixed. If configuration $b$ is reachable from configuration $a$ by a finite run whose
configurations before the endpoint satisfy $G$, then any configuration equal to $b$ is
also reachable from $a$ by such a run.
\end{lemma}

\begin{theorem}[Register Transfer Loop]
\lean{Complexity.LogProg.move_run}
\label{Complexity.LogProg.move_run}
\leanok
\uses{Complexity.LogProg.AReach, Complexity.LogProg.AReach.refl, Complexity.LogProg.AReach.step, Complexity.LogProg.GuardLe, Complexity.LogProg.PreS, Complexity.LogProg.astep, Complexity.LogProg.reach_of_eq}
\difficulty{5}
Let $A$ be a register program with $m$ registers and $d$ oracle deciders, let $x$ be an
input, and let $B$ be a natural-number bound. Suppose distinct registers $r,t$ and
states $l,l_x,l_1,l_2$ satisfy: at $l$ the program branches to $l_x$ when register $r$
is zero and to $l_1$ otherwise; at $l_1$ it decrements $r$ and proceeds to $l_2$; and at
$l_2$ it increments $t$ and returns to $l$. For every $k$ and register assignment $v$
with $v(r)=k$, with every register value at most $B$, and with $v(t)+v(r)\le B$, the
configuration at $l$ reaches the configuration at $l_x$ with $r$ set to zero and $t$ set
to its original value $v(t)+v(r)$, all other registers unchanged, and no answer
recorded; every configuration before the final one satisfies the program's precondition
and has all register values at most $B$.
\end{theorem}

\begin{theorem}[Bounded run for a register transfer]
\lean{Complexity.LogProg.move2_run}
\label{Complexity.LogProg.move2_run}
\leanok
\uses{Complexity.LogProg.AReach, Complexity.LogProg.AReach.refl, Complexity.LogProg.AReach.step, Complexity.LogProg.GuardLe, Complexity.LogProg.PreS, Complexity.LogProg.astep, Complexity.LogProg.reach_of_eq, Complexity.LogProg.register_update_le}
\difficulty{5}
Let $A$ be a register program, let an oracle and input word be fixed, and let $B$ be a
natural-number bound. Suppose distinct registers $r,t,t'$ and states $l,l_x,l_1,l_2,l_3$
satisfy: at $l$ the program branches on whether register $r$ is zero, sending the zero
case to $l_x$ and the other case to $l_1$; at $l_1$ it decrements $r$ and continues to
$l_2$; at $l_2$ it increments $t$ and continues to $l_3$; and at $l_3$ it increments
$t'$ and returns to $l$. For every natural number $k$ and register valuation $v$ with
$v(r)=k$, all register values at most $B$, and $v(t)+v(r)\le B$ and $v(t')+v(r)\le B$,
there is a finite run from $(l,v,\mathrm{none})$ to $(l_x,v',\mathrm{none})$ in which
$r$ is set to $0$, while $t$ and $t'$ are each increased by the initial value $v(r)$ and
every other register retains its value. At every configuration before the final one, the
active instruction satisfies its applicable side condition (register comparisons use
distinct registers, input comparisons use valid encodings, and call arguments are
distinct), and every register value is at most $B$.
\end{theorem}

\begin{theorem}[Bounded register transfer run]
\lean{Complexity.LogProg.dmove_run}
\label{Complexity.LogProg.dmove_run}
\leanok
\uses{Complexity.LogProg.AReach, Complexity.LogProg.AReach.refl, Complexity.LogProg.AReach.step, Complexity.LogProg.GuardLe, Complexity.LogProg.PreS, Complexity.LogProg.astep, Complexity.LogProg.reach_of_eq, Complexity.LogProg.register_update_le}
\difficulty{5}
Let $A$ be a register program, fix an oracle and input, and suppose registers $r$ and
$t$ are distinct. Suppose the instructions at states $l,l_1,l_2,l_3$ are, respectively,
a zero test on $r$ branching to $l_x$ or $l_1$, a decrement of $r$ leading to $l_2$, an
increment of $t$ leading to $l_3$, and an increment of $t$ leading back to $l$. For
every $k$ and register valuation $v$ with $v(r)=k$, all registers bounded above by $B$,
and $v(t)+2v(r)\le B$, there is a finite run from state $l$ with valuation $v$ to state
$l_x$ with $r$ reset to zero and $t$ increased by $2v(r)$, leaving the other registers
unchanged; every configuration before the endpoint satisfies the guard that its
precondition holds and all registers are at most $B$.
\end{theorem}

\begin{theorem}[Register Subroutine Run]
\lean{Complexity.LogProg.sub_run}
\label{Complexity.LogProg.sub_run}
\leanok
\uses{Complexity.LogProg.AReach, Complexity.LogProg.AReach.refl, Complexity.LogProg.AReach.step, Complexity.LogProg.GuardLe, Complexity.LogProg.PreS, Complexity.LogProg.astep, Complexity.LogProg.reach_of_eq, Complexity.LogProg.register_update_le}
\difficulty{5}
Let $A$ be a register program, let $q,r$ be distinct registers, and suppose its
instructions at states $l,l_1,l_2$ are respectively: branch to $l_x$ when register $q$
is zero (with the other branch to $l_1$), decrement $q$ and go to $l_2$, then decrement
$r$ and return to $l$. For every $k\in\mathbb N$ and register valuation $v$ with
$v(q)=k$ and every register value at most $B$, there is a finite run from state $l$ to
state $l_x$ that satisfies the guard requiring the program's side conditions and every
register value to be at most $B$ at each earlier step. The final valuation sets $q$ to
$0$, decreases $r$ by $k$ using truncated subtraction, and leaves all other registers
unchanged; the final result is unset.
\end{theorem}

\begin{theorem}[Repeated halving after a countdown]
\lean{Complexity.LogProg.halve_run}
\label{Complexity.LogProg.halve_run}
\leanok
\uses{Complexity.LogProg.AReach, Complexity.LogProg.AReach.refl, Complexity.LogProg.AReach.step, Complexity.LogProg.GuardLe, Complexity.LogProg.PreS, Complexity.LogProg.astep, Complexity.LogProg.reach_of_eq, Complexity.LogProg.register_update_le}
\difficulty{6}
Let $A$ be a register program, let $l,l_x,l_1,l_2$ be labels, and let $c,r$ be distinct
registers. Suppose label $l$ tests whether register $c$ is zero, its positive branch
decrements $c$ and goes to $l_2$, and $l_2$ halves register $r$ and returns to $l$. For
every $k\in\mathbb N$ and register assignment $v$ with $v(c)=k$ and $v(s)\le B$ for
every register $s$, the program can reach the configuration at $l_x$ obtained by setting
$c$ to zero and replacing $r$ by $v(r)/2^{v(c)}$, with no result bit, while every
configuration before the final one satisfies the guard requiring the program's
precondition and all registers to be at most $B$.
\end{theorem}

\begin{theorem}[Bounded run of a register copy routine]
\lean{Complexity.LogProg.copy_run}
\label{Complexity.LogProg.copy_run}
\leanok
\uses{Complexity.LogProg.AReach, Complexity.LogProg.AReach.trans, Complexity.LogProg.GuardLe, Complexity.LogProg.move2_run, Complexity.LogProg.move_run, Complexity.LogProg.reach_of_eq, Complexity.LogProg.register_update_le}
\difficulty{4}
Let $A$ be a register program with oracle and let $B$ be a natural number. Choose
distinct registers $src,dst,z$ and states $l,l_1,l_2,l_3,l_4,l_5,l_6,l_x$. Suppose the
program instructions at these states implement the following routine: at $l$, branch on
whether $src=0$; if nonzero, decrement $src$, increment $dst$, and increment $z$, then
return to $l$; if zero, branch on whether $z=0$; while $z$ is nonzero, decrement $z$,
increment $src$, and repeat that branch; when $z=0$, proceed to $l_x$. For every initial
register assignment $v$ with $v(z)=0$, all register values at most $B$, and
$v(dst)+v(src)\le B$, there is a finite run from state $l$ to state $l_x$ that changes
only $dst$, setting it to $v(dst)+v(src)$, and leaves every register value at most $B$.
Every configuration before the final one satisfies the program's guard: any
equality-branch registers are distinct, any input-dependent checks are valid on $x$, and
any call uses an argument list without repetitions.
\end{theorem}

\begin{theorem}[Register doubling by a logarithmic loop]
\lean{Complexity.LogProg.shl_run}
\label{Complexity.LogProg.shl_run}
\leanok
\uses{Complexity.LogProg.AReach, Complexity.LogProg.AReach.refl, Complexity.LogProg.AReach.step, Complexity.LogProg.AReach.trans, Complexity.LogProg.GuardLe, Complexity.LogProg.PreS, Complexity.LogProg.astep, Complexity.LogProg.dmove_run, Complexity.LogProg.move_run, Complexity.LogProg.reach_of_eq, Complexity.LogProg.register_update_le}
\difficulty{6}
Let $A$ be an instruction program on $m$ registers, with configurations consisting of a
control state, register values in $\bbn$, and an optional Boolean result; let an oracle
assign a Boolean answer to each of $d$ deciders on Boolean-string inputs. Suppose the
transitions from states $l,l_1,\ldots,l_8$ are respectively: test whether register $c$
is zero (branching to $l_x$ if so and $l_1$ otherwise); decrement $c$ and go to $l_2$;
test whether $r$ is zero (branching to $l_6$ if so and $l_3$ otherwise); decrement $r$
and go to $l_4$; increment $z$ and go to $l_5$; increment $z$ and return to $l_2$; test
whether $z$ is zero (branching to $l$ if so and $l_7$ otherwise); decrement $z$ and go
to $l_8$; and increment $r$ and return to $l_6$. Assume $c,r,z$ are pairwise distinct
registers. For every $k	rin\bbn$ and register assignment $v$ with $v(c)=k$, $v(z)=0$,
every register value at most $B$, and $v(r)2^{v(c)}	rin B$, this program reaches the
configuration with state $l_x$, register $c$ set to zero, register $r$ set to
$v(r)2^{v(c)}$, all other registers unchanged, and no result. Every configuration
strictly before that endpoint satisfies the guard: it meets the instruction-specific
safety condition and all register values are at most $B$.
\end{theorem}
